\chapter{Logarithms and Exponentials}

Mix of evaluating, expanding, condensing, and solving. Bases of 10 and 2
appear alongside the natural log.

\section{Core Rules}

\begin{rulebox}{What a logarithm is}
$\log_{a}(b)$ is \textbf{the exponent that turns $a$ into $b$}:
\[
  \log_{a}(b)=c\iff a^{c}=b\qquad (a>0,\ a\neq1,\ b>0).
\]
Every log identity below is an exponent law read backwards. The domain
restriction $b>0$ is part of the definition, not a technicality: it is what
manufactures extraneous roots when you solve equations.
\end{rulebox}

\begin{rulebox}{The three laws, and the change of base}
\[
  \ln(xy)=\ln x+\ln y,\qquad
  \ln\!\Bigl(\frac{x}{y}\Bigr)=\ln x-\ln y,\qquad
  \ln(x^{p})=p\ln x .
\]
The third is the workhorse: it is what pulls an unknown down out of an
exponent. Any base converts to any other,
\[
  \log_{a}(b)=\frac{\ln b}{\ln a},
\]
so a calculator with only $\ln$ suffices. Also $\ln(\mathrm{e}^{x})=x$ and
$\mathrm{e}^{\ln x}=x$: the two functions undo each other.
\end{rulebox}

\begin{rulebox}{There is no law for a sum inside the log}
\[
  \ln(x+y)\neq\ln x+\ln y .
\]
The left side cannot be expanded at all. The laws convert
\emph{multiplication into addition}, never addition into anything. The same
warning in the other direction: $\ln x\cdot\ln y$ does not simplify, and
$(\ln x)^{2}\neq\ln(x^{2})=2\ln x$.
\end{rulebox}

\begin{rulebox}{Solving: the standard shapes}
\begin{itemize}[leftmargin=1.6em,itemsep=0.3ex]
  \item \textbf{Unknown in the exponent} ($a^{x}=b$): take $\ln$ of both
        sides, then $x=\ln b/\ln a$.
  \item \textbf{Unknown in the base} ($x^{n}=b$): raise both sides to
        $1/n$.
  \item \textbf{Logs on both sides}: condense each side to a single log,
        then drop the logs --- legitimate because $\ln$ is one-to-one.
  \item \textbf{One log, one number}: condense, then exponentiate.
\end{itemize}
Then \textbf{check every root against the original domain.}
\end{rulebox}

\section{Worked Examples}

\begin{worked}{(a) Solve $\log_{3}(x)+\log_{3}(x-2)=1$}
\wstep{Step 1. Record the domain first, before solving anything.}
Both logs must take positive arguments: $x>0$ \emph{and} $x-2>0$, so
\[
  x>2 .
\]
Writing this down now costs one line and settles the last step
automatically.

\wstep{Step 2. Condense the left side to a single logarithm.}
\[
  \log_{3}(x)+\log_{3}(x-2)=\log_{3}\bigl(x(x-2)\bigr).
\]

\wstep{Step 3. Undo the logarithm.}
The equation now reads $\log_{3}\bigl(x(x-2)\bigr)=1$, and by the
definition of a logarithm that says $3^{1}=x(x-2)$:
\[
  x^{2}-2x=3\quad\Longrightarrow\quad x^{2}-2x-3=0
  \quad\Longrightarrow\quad (x-3)(x+1)=0 .
\]

\wstep{Step 4. Apply the domain.}
The candidates are $x=3$ and $x=-1$. The second fails $x>2$ --- indeed
$\log_{3}(-1)$ does not exist --- so
\[
  \boxed{x=3}.
\]
Check: $\log_{3}3+\log_{3}1=1+0=1$. \checkmark
\wtrap{$x=-1$ is not a wrong calculation; it is a genuine root of the
\emph{quadratic}, which is a different equation from the one asked. The
condensing step in Step~2 enlarged the domain, and enlarging a domain is
exactly what creates extraneous roots. Any time you condense two logs into
one, the final domain check is compulsory.}
\end{worked}

\begin{worked}{(b) How long until an economy doubles? Solve $Y_{0}\mathrm{e}^{gt}=2Y_{0}$, and solve $(1+r)^{10}=2$}
\wstep{Step 1. Cancel what is common, then isolate the exponential.}
$Y_{0}>0$, so dividing by it is legal and the initial level drops out
entirely: $\mathrm{e}^{gt}=2$. \emph{Doubling time does not depend on where
you start} --- that is the content of exponential growth, and it is visible
here before any calculation.

\wstep{Step 2. The unknown is in the exponent, so take logs.}
\[
  \ln\bigl(\mathrm{e}^{gt}\bigr)=\ln 2
  \quad\Longrightarrow\quad
  gt=\ln 2
  \quad\Longrightarrow\quad
  t=\frac{\ln 2}{g}\approx\frac{0.693}{g}.
\]
At $g=0.03$ this is $23.1$ years --- the ``rule of 70'', which is just
$100\ln 2\approx69.3$ rounded.

\wstep{Step 3. The second equation has the unknown in the \emph{base}.}
$(1+r)^{10}=2$: here $r$ sits inside the base, not the exponent, so logs
are the wrong tool. Undo the tenth power with a tenth root:
\[
  1+r=2^{1/10}\quad\Longrightarrow\quad r=2^{1/10}-1\approx0.0718 .
\]
So $7.18\%$ a year doubles capital in a decade.
\wtrap{Decide \emph{where the unknown lives} before choosing a move.
Unknown in the exponent $\Rightarrow$ take logs. Unknown in the base
$\Rightarrow$ take a root. Taking logs of $(1+r)^{10}=2$ is not wrong, but
it leads to $\ln(1+r)=\tfrac{1}{10}\ln2$ and one more exponentiation to
undo --- the long way round.}
\end{worked}

\section{Practice Problems}

\begin{problems}
  \item Evaluate: $\ln(1)$
  \item Evaluate: $\log_{10}(100)$
  \item Evaluate: $\log_{2}(32)$
  \item Evaluate: $\ln(\mathrm{e})+\ln(\mathrm{e}^{3})$
  \item Expand: $\displaystyle\ln\left(\frac{x^{a}}{y^{b}z^{c}}\right)$
  \item Expand: $\displaystyle\ln\left(\sqrt[5]{\frac{x^{3}}{y^{2}}}\right)$
  \item Can $\ln(x^{2}+y^{2})$ be expanded? Why?
  \item Expand: $\ln(C\,\mathrm{e}^{-\delta t})$
  \item Condense: $4\ln(x)-2\ln(y)+\tfrac12\ln(z)$
  \item Condense: $\tfrac13\ln(8)+\tfrac12\ln(9)$
  \item Condense: $\ln(a)+\ln(b)-\ln(c)-\ln(d)$
  \item Solve for $x$: $\log_{10}(x)=3$
  \item Solve for $x$ ($x>0$): $\ln(x)=-1$
  \item Solve for $x$ ($x>2$): $\ln(x)+\ln(x-2)=\ln(3)$
  \item Half-life: solve $\tfrac12 P_{0}=P_{0}\,\mathrm{e}^{-\lambda t}$ for
        $t$.
  \item Solve for $r$: $(1+r)^{10}=2$
  \item Solve for $\alpha$ ($0<\alpha<1$): $\ln(\alpha)-\ln(1-\alpha)=1$
  \item Solve for $x$: $\mathrm{e}^{2x}-3\mathrm{e}^{x}+2=0$ (Hint:
        substitute $u=\mathrm{e}^{x}$.)
  \item Solve for $T$: $A=P(1+r)^{T}$
  \item Solve for $x$ ($x>0$): $(\ln x)^{2}=\ln(x^{4})$
  \item Solve for $t$: $A_{0}\cdot 2^{-t/h}=\dfrac{A_{0}}{8}$
  \item Expand: $\displaystyle\ln\left(\frac{(1+r)^{T}}{r}\right)$
  \item Solve: $\log_{3}(x)+\log_{3}(x-2)=1$
  \item Solve: $\log_{2}(x+1)-\log_{2}(x-1)=1$
  \item Solve: $\mathrm{e}^{x}=5\,\mathrm{e}^{-x}$
  \item Solve: $4^{x}=8$
  \item Given $\ln 2\approx 0.693$ and $\ln 3\approx 1.099$, estimate
        $\ln 12$.
  \item Use the change-of-base formula to express $\log_{5}(20)$ as a ratio
        of natural logs.
  \item Solve for $K$ in $\ln(Y)=\ln(A)+\alpha\ln(K)+(1-\alpha)\ln(L)$.
  \item Solve for $g$: $Y_{t}=Y_{0}\,\mathrm{e}^{gt}$
\end{problems}

\section{Solutions}

\begin{solutions}
  \item $0$
  \item $2$
  \item $5$
  \item $1+3=4$
  \item $a\ln(x)-b\ln(y)-c\ln(z)$
  \item $\tfrac15\bigl[3\ln(x)-2\ln(y)\bigr]$
  \item No: $\ln$ has no product/quotient rule for a sum inside.
  \item $\ln(C)-\delta t$
  \item $\displaystyle\ln\left(\frac{x^{4}\sqrt{z}}{y^{2}}\right)$
  \item $\ln(2)+\ln(3)=\ln(6)$
  \item $\displaystyle\ln\left(\frac{ab}{cd}\right)$
  \item $x=1000$
  \item $x=\dfrac{1}{\mathrm{e}}$
  \item $x(x-2)=3\Rightarrow x=3$
  \item $t=\dfrac{\ln 2}{\lambda}$
  \item $r=2^{1/10}-1$
  \item $\displaystyle\frac{\alpha}{1-\alpha}=\mathrm{e}\Rightarrow
        \alpha=\frac{\mathrm{e}}{1+\mathrm{e}}$
  \item $\mathrm{e}^{x}=1$ or $\mathrm{e}^{x}=2\Rightarrow x=0$ or
        $x=\ln 2$
  \item $T=\dfrac{\ln(A/P)}{\ln(1+r)}$
  \item $\ln x=0$ or $\ln x=4\Rightarrow x=1$ or $x=\mathrm{e}^{4}$
  \item $2^{-t/h}=2^{-3}\Rightarrow t=3h$
  \item $T\ln(1+r)-\ln(r)$
  \item $x(x-2)=3\Rightarrow x=3$
  \item $\dfrac{x+1}{x-1}=2\Rightarrow x=3$
  \item $\mathrm{e}^{2x}=5\Rightarrow x=\tfrac12\ln 5$
  \item $x=\tfrac32$
  \item $\ln 12=2\ln 2+\ln 3\approx 2.485$
  \item $\log_{5}(20)=\dfrac{\ln 20}{\ln 5}$
  \item $K=\displaystyle\left(\frac{Y}{AL^{1-\alpha}}\right)^{1/\alpha}$
  \item $g=\dfrac{1}{t}\ln\!\left(\dfrac{Y_{t}}{Y_{0}}\right)$
\end{solutions}
